\section{Structure des programmes} 
%\label{sec:4:ecrir-progs} 

De même que pour écrire correctement en français il faut respecter
des règles précises d'orthographe et de grammaire, l'écriture de
programmes corrects est soumise à un certain nombre de {\em règles
de syntaxe}\index{Règle de syntaxe} heureusement plus simples que pour le français. La différence
est que des phrases françaises comportant des fautes d'orthographe ou
de grammaire peuvent rester compréhensibles par des francophones alors
qu'un programme comportant des 
{\em erreurs syntaxiques}\index{Erreur syntaxique} ne peut
pas être traduit dans le langage de l'or\-di\-na\-teur, ni donc exécuté.
Cette traduction est effectuée par un programme (appelé {\em interpréteur})\index{Interpréteur}
prévu pour indiquer nos erreurs mais pas pour les corriger. Pour corriger
nos erreurs de manière satisfaisante, l'interpréteur devrait pouvoir comprendre 
nos intentions, un problème bien difficile.

Le programme {\idt Cocotte} ci-dessous est un exemple de
programme \python\ correct.

\begin{exemple}[Programme syntaxiquement correct]{\hfill
\label{prog:Cocotte} \input python/Cocotte.tex
\noindent $\Box$ }\end{exemple}

\noindent Tous nos programmes \python\ de pilotage du \robot\ ont la même forme :

\begin{center}
\begin{tabular}{@{}|p{12.5cm}|@{}} \hline
{\em Programme \python\ :}\\ \hline \hline
\begin{minipage}[t]{12.5cm}
{\medskip
\mbox{\blc\ \ \ {\mcl from}\hspace{0.3em}\blanc\hspace{0.3em}{\idt LibRobot}\hspace{0.3em}\blanc\hspace{0.3em}{\mcl import}\hspace{0.3em}\blanc\hspace{0.3em}{\idt *}}\par 
\mbox{\blc\ \ \ \cmt \#\ \small\rm Commentaire expliquant ce que fait le programme}\par 
\mbox{\blc\ \ \ \small\rm  \ldots\ Instructions du programme (une par ligne, ou séparées par des `$;$') \ldots}\par 
\medskip}
\end{minipage}\\ \hline 
\end{tabular}
\end{center}

Un programme \python\ est rangé dans un fichier dont le nom se termine par {\idt .py}.
Le nom du programme est celui de ce fichier, sans le {\idt .py}~: le programme 
{\idt Cocotte} est rangé dans le fichier {\idt Cocotte.py}.

Ce nom est choisi par le programmeur
à sa convenance, de façon à évoquer ce à quoi sert son programme.
 La seule règle à res\-pec\-ter est que le nom du programme
ne doit comporter que des lettres minuscules ou
majuscules non accentuées\footnote{\ Les accents n'existent pas en anglais
(à de très rares exceptions près, comme dans les mots `naïve', `protégé' ou
`rôle' empruntés au français).}, des chiffres ou
des soulignés `\_'\footnote{\ Le souligné `\_' ne doit pas être confondu avec le tiret `-'.}
et doit commencer par une lettre. Par exemple {\idt  S}, 
{\idt  LeStegosaure} ou {\idt  Le\_Stegosaure} sont des noms corrects de
programmes alors que {\idt  1S}, 
{\idt  s!}, {\idt  Le\ Stegosaure} ou {\idt  Le-Stégosaure} ne respectent
pas la règle d'écriture des noms de programmes. 

Un programme \python\ commence par la ligne `{\mcl from}\ {\idt LibRobot}\ {\mcl import}\ {\idt *}'.
Les mots anglais {\mcl from} et {\mcl import} signifient `de' et `importer'~: cette ligne 
importe de la 
{\em bibliothèque}\index{Bibliothèque} {\idt LibRobot} toutes (c'est le sens de l'étoile `{\idt *}') 
les commandes du 
\robot, sans quoi l'ordinateur ne les connaîtrait pas. Elle ouvre aussi la fenêtre 
dans laquelle le \robot\ dessine.

Il y a obligatoirement un blanc
entre le mot {\mcl from} et le nom {\idt LibRobot}, ainsi qu'entre {\idt LibRobot} et le 
mot {\mcl import}. Dans un livre il y a des blancs partout et même, en y regardant bien,
plus de blanc que toute autre chose. Mettre un blanc dans un livre pour dire qu'il faut taper
un blanc est donc idiot parce que la plupart des lecteurs ne s'intéressent pas aux 
blancs dans les livres mais simplement à ce qui y est écrit.
Nous utiliserons donc le symbole `\blanc' pour écrire nos blancs noir sur \hbox{blanc$\:$!}
Autrement dit nous écrivons \blanc\ pour dire qu'il faut taper sur la grande barre 
située en bas du clavier de l'or\-di\-na\-teur qui s'appelle la 
{\em barre d'espacement}\index{Barre d'espacement}
pour obtenir un blanc `\ ' dans le programme. 
Nous ne le faisons pas dans les programmes donnés en exemple qui sont
présentés comme ils apparaissent à l'écran. De même, à la fin d'une 
ligne il faut taper sur la touche de retour à la ligne 
\framebox[11.mm][l]{$\hookleftarrow$~~~}\hspace{0.3ex} ou \hspace{0.3ex}\framebox[11.mm][l]{~~\tiny\texttt{return}}\hspace{0.3ex} pour passer au début de la
ligne suivante. Nous ne l'écrivons pas dans les livres car tous les lecteurs savent
qu'il faut passer au début de la ligne suivante quand ils sont arrivés en fin de ligne. 

Après le commentaire qui explique ce que fait le programme
viennent les instructions 
du programme. Elles sont écrites dans l'ordre où elles doivent être exécutées par
l'or\-di\-na\-teur, en commençant au début de la ligne (sans blanc devant~: nous verrons 
au chapitre suivant que les blancs en début de ligne ont un sens précis en \python). 
On peut écrire une instruction par ligne ou bien plusieurs instructions sur une même ligne 
en les séparant par des points-virgules, sans point-virgule après la dernière de la ligne. 
Les instructions du programme
peuvent être des commandes du \hbox{\robot\ :} {\idt av}, {\idt vg}, {\idt vd}, {\idt pqtd}, 
{\idt pqtg}, {\idt phtd}, {\idt phtg}, {\idt lc}, {\idt bc}, {\idt dc}, {\idt dp}, 
{\idt dg} ou {\idt st}, chacune suivie de {\idt ()}. En général les programmes de pilotage du \robot\ se
terminent par la commande {\idt st()} qui permet d'attendre que nous
tapions sur une touche avant que le dessin ne 
disparaisse à la fin de l'exécution du programme.

Les textes qui commencent par le signe {\cmt \#} et se poursuivent jusqu'à la fin de 
la ligne sont des {\em commentaires} ignorés par l'ordinateur.

Le pro\-gram\-me
{\idt CocotteErronee} ci-dessous contient
plusieurs erreurs très courantes chez les débutants.

\begin{exemple}[Programme syntaxiquement incorrect]{\hfill

\label{prog:Cocotte-errone} \input python/CocotteErronee.tex
\noindent Les erreurs dans ce programme sont les \hbox{suivantes :}
\begin{list}{$\scriptscriptstyle\bullet$}{
\renewcommand{\labelwidth}{3mm}
\renewcommand{\labelsep}{3mm}
\renewcommand{\leftmargin}{10mm}
\renewcommand{\rightmargin}{0mm}
}
\item[1.] la ligne `{\mcl from}\ {\idt LibRobot}\ {\mcl import}\ {\idt *}' a été oubliée$\:$;
\item[2.] le texte qui explique ce que fait le programme a été écrit entre triples 
guillemets {\cha \char34\char34\char34} (c'est une {\em chaîne de documentation}, que l'on 
peut utiliser comme un commentaire de plusieurs lignes) mais les triples guillemets 
fermants {\cha \char34\char34\char34} ont été oubliés$\:$;
\item[3.] l'instruction {\idt pgtg()} n'est pas une commande du \robotErr$\:$;
\item[4.] la parenthèse ouverte de la dernière commande {\idt st(} n'est pas fermée.
\end{list}
\noindent $\Box$ }\end{exemple}


Un exemple compliqué de programme \python\ syntaxiquement correct
est donné ci-\hbox{dessous :}
\begin{exemple}[Le stégosaure]{\rm
Le programme \python\ de pilotage du \robot\ pour
dessiner le sté\-go\-sau\-re re\-pré\-sen\-té  à la fi\-gu\-re
\ref{fig:stegosaure} (page \pageref{fig:stegosaure}) est le \hbox{suivant :}
\label{prog:Stegosaure} \input python/Stegosaure.tex
\noindent $\Box$ }\end{exemple}

\section{Commentaires} 
%\label{sec:4:comment} 

Un programme \python\ peut comporter des {\em commentaires}\index{Commentaire} qui sont
des textes français quelconques écrits après le signe {\cmt \#\ \ldots}, jusqu'à la fin de la 
ligne. Un commentaire peut occuper toute une ligne ou bien suivre les instructions 
d'une ligne, comme dans {\idt vg()\ \ }{\cmt \#\ ERREUR\ :\ vd}.
Un texte de plusieurs lignes peut aussi être placé entre triples guillemets 
{\cha \char34\char34\char34\ \ldots\ \char34\char34\char34}~: c'est une 
{\em chaîne de documentation}\index{Chaîne de documentation}, que l'ordinateur 
lit mais qui ne commande rien au \robot.
L'or\-di\-na\-teur ne tient pas compte de ces commentaires qui sont uniquement
destinés à expliquer l'usage et le principe de fonctionnement du
programme aux lecteurs de ce programme. Si l'on enlève tous les
commentaires d'un programme, il devient illisible pour un humain mais
rien n'est changé pour l'or\-di\-na\-teur.

\begin{exemple}[Le stégosaure \underline{illisible}]{\rm
En enlevant tous les com\-men\-tai\-res du programme {\idt Stegosaure}
de la page \pageref{prog:Stegosaure}, ainsi que les blancs et les passages à la 
ligne inutiles, on obtient un programme \python\
totalement illisible mais qui dessine parfaitement bien le 
sté\-go\-sau\-re re\-pré\-sen\-té  à la fi\-gu\-re
\ref{fig:stegosaure} (page \hbox{\pageref{fig:stegosaure}) :}
{\small{
\parindent \parindentprog
\medskip

\carrenoir
\vspace{\neglskip}\mbox{\blc {\mcl from}{\idt \ LibRobot\ }{\mcl import}{\idt \ *}}\par
\vspace{\neglskip}\mbox{\blc {\idt lc();\ pqtd();\ pqtd();\ av();\ av();\ av();\ av();\ av();\ av();\ av()}}\par
\vspace{\neglskip}\mbox{\blc {\idt pqtg();\ av();\ av();\ av();\ av();\ av();\ av();\ av();\ av();\ av()}}\par
\vspace{\neglskip}\mbox{\blc {\idt av();\ av();\ av();\ av();\ pqtg();\ bc();\ av();\ pqtd();\ vg();\ pqtg()}}\par
\vspace{\neglskip}\mbox{\blc {\idt vg();\ pqtd();\ pqtd();\ vg();\ pqtd();\ pqtd();\ vg();\ pqtg();\ vg()}}\par
\vspace{\neglskip}\mbox{\blc {\idt pqtd();\ vd();\ vg();\ av();\ vd();\ av();\ av();\ av();\ pqtd();\ av()}}\par
\vspace{\neglskip}\mbox{\blc {\idt pqtd();\ av();\ av();\ av();\ av();\ pqtd();\ vd();\ av();\ pqtg();\ vd()}}\par
\vspace{\neglskip}\mbox{\blc {\idt av();\ av();\ av();\ pqtd();\ av();\ pqtd();\ av();\ av();\ pqtd();\ lc()}}\par
\vspace{\neglskip}\mbox{\blc {\idt av();\ bc();\ vd();\ av();\ av();\ pqtd();\ av();\ pqtd();\ av();\ pqtd()}}\par
\vspace{\neglskip}\mbox{\blc {\idt lc();\ av();\ bc();\ vd();\ pqtg();\ av();\ pqtd();\ av();\ pqtd();\ av()}}\par
\vspace{\neglskip}\mbox{\blc {\idt av();\ pqtd();\ av();\ av();\ pqtd();\ vd();\ pqtg();\ av();\ av();\ vd()}}\par
\vspace{\neglskip}\mbox{\blc {\idt av();\ pqtd();\ av();\ av();\ pqtd();\ av();\ av();\ pqtd();\ av();\ vd()}}\par
\vspace{\neglskip}\mbox{\blc {\idt pqtg();\ av();\ av();\ pqtd();\ av();\ av();\ pqtd();\ av();\ av()}}\par
\vspace{\neglskip}\mbox{\blc {\idt pqtd();\ av();\ av();\ pqtd();\ av();\ av();\ pqtg();\ vd();\ av()}}\par
\vspace{\neglskip}\mbox{\blc {\idt pqtd();\ av();\ av();\ pqtd();\ av();\ av();\ pqtd();\ av();\ vd();\ av()}}\par
\vspace{\neglskip}\mbox{\blc {\idt av();\ pqtg();\ vd();\ pqtd();\ av();\ av();\ pqtd();\ av();\ av()}}\par
\vspace{\neglskip}\mbox{\blc {\idt pqtd();\ av();\ pqtd();\ av();\ pqtg();\ vd();\ vg();\ pqtd();\ pqtd()}}\par
\vspace{\neglskip}\mbox{\blc {\idt av();\ av();\ pqtd();\ av();\ pqtd();\ av();\ pqtd();\ vg();\ vd();\ vg()}}\par
\vspace{\neglskip}\mbox{\blc {\idt pqtd();\ pqtd();\ av();\ av();\ av();\ pqtd();\ av();\ av();\ pqtd()}}\par
\vspace{\neglskip}\mbox{\blc {\idt av();\ vd();\ vg();\ av();\ vd();\ pqtd();\ av();\ av();\ av();\ pqtd()}}\par
\vspace{\neglskip}\mbox{\blc {\idt av();\ pqtd();\ av();\ av();\ vd();\ av();\ av();\ lc();\ pqtg();\ av()}}\par
\vspace{\neglskip}\mbox{\blc {\idt dc();\ pqtd();\ pqtd();\ av();\ pqtg();\ bc();\ vg();\ vd();\ av();\ vg()}}\par
\vspace{\neglskip}\mbox{\blc {\idt pqtg();\ av();\ av();\ pqtd();\ pqtd();\ av();\ av();\ pqtg();\ vg()}}\par
\vspace{\neglskip}\mbox{\blc {\idt av();\ av();\ av();\ av();\ av();\ av();\ av();\ vg();\ pqtd();\ av()}}\par
\vspace{\neglskip}\mbox{\blc {\idt pqtg();\ av();\ av();\ av();\ pqtd();\ pqtd();\ av();\ av();\ av();\ av()}}\par
\vspace{\neglskip}\mbox{\blc {\idt av();\ pqtg();\ av();\ av();\ pqtg();\ av();\ av();\ av();\ av();\ av()}}\par
\vspace{\neglskip}\mbox{\blc {\idt av();\ pqtd();\ pqtd();\ av();\ av();\ av();\ av();\ av();\ pqtg();\ av()}}\par
\vspace{\neglskip}\mbox{\blc {\idt pqtg();\ av();\ av();\ pqtg();\ av();\ pqtd();\ pqtd();\ av();\ vg()}}\par
\vspace{\neglskip}\mbox{\blc {\idt pqtd();\ av();\ av();\ av();\ pqtg();\ av();\ av();\ av();\ pqtd()}}\par
\vspace{\neglskip}\mbox{\blc {\idt pqtd();\ av();\ av();\ av();\ av();\ av();\ av();\ av();\ pqtg();\ av()}}\par
\vspace{\neglskip}\mbox{\blc {\idt av();\ pqtg();\ av();\ av();\ av();\ av();\ av();\ av();\ av();\ av()}}\par
\vspace{\neglskip}\mbox{\blc {\idt pqtd();\ pqtd();\ av();\ av();\ av();\ av();\ av();\ av();\ av();\ pqtg()}}\par
\vspace{\neglskip}\mbox{\blc {\idt av();\ pqtg();\ av();\ av();\ av();\ pqtg();\ av();\ pqtd();\ pqtd()}}\par
\vspace{\neglskip}\mbox{\blc {\idt av();\ av();\ pqtd();\ vg();\ av();\ vd();\ vg();\ av();\ pqtd();\ vg()}}\par
\vspace{\neglskip}\mbox{\blc {\idt pqtd();\ st()}}\par

\medskip }
}
Cet exemple montre également que la disposition des instructions du
programme est importante. L'espacement entre instructions et les
alignements convenables facilitent grandement la lecture des
programmes.
\noindent $\Box$ }\end{exemple}



\section{Interprétation des programmes} 
%\label{sec:4:compil-progs} 
L'{\em interprétation}\index{Interprétation} d'un programme \python\ consiste à le traduire 
dans le langage de l'or\-di\-na\-teur
appelé {\em langage machine}\index{Langage machine} et à le faire exécuter par l'ordinateur. 
Elle est faite par un programme
que l'on appelle l'\index{Interpréteur}{\em interpréteur} \python\@\footnote{\ La commande 
{\tt ./robot.sh Cocotte.py}, tapée dans une 
fenêtre de terminal placée dans le dossier {\tt LeRobot\_Python\_FR}, appelle l'interpréteur 
{\tt python3} pour faire exécuter le programme {\tt Cocotte.py}. Faire glisser l'icône de 
{\tt Cocotte.py} sur celle de l'application {\tt LeRobotPython.app} fait la même chose.}.
L'interpréteur procède en deux temps~:
\begin{list}{$\scriptscriptstyle\bullet$}{
\renewcommand{\rightmargin}{0mm}
}
\item il lit d'abord {\em tout} le programme pour vérifier qu'il respecte les règles 
de syntaxe de \python\ et le traduire$\:$; s'il trouve une erreur syntaxique, il le signale 
et s'arrête sans rien exécuter$\:$;
\item puis, si le programme est syntaxiquement correct, il fait exécuter ses 
instructions par l'ordinateur, les unes après les autres.
\end{list}
Certaines erreurs ne peuvent être découvertes que pendant ce second temps, au moment où 
l'or\-di\-na\-teur arrive à l'instruction fautive, par exemple l'appel d'une commande qui n'existe 
pas. L'exécution s'arrête alors, mais le \robot\ a déjà fait tout le dessin demandé 
par les instructions précédentes.
Bien évidemment le programme \python\ et sa traduction pilotent
le \robot\ de la même façon. Cette manière de procéder nous 
évite d'avoir à apprendre le langage machine. Il est bien difficile
de comprendre ce langage machine car tous les programmes s'écrivent
uniquement avec des 0 et \hbox{des 1$\:$!}

\section{Erreurs syntaxiques} 
%\label{sec:4:err-compil} 

Pour pouvoir exécuter un programme \python, il ne faut pas qu'il contienne
d'erreurs syntaxiques\index{Erreur de programmation}\index{Erreur syntaxique} violant les règles
d'écriture des programmes. En cas d'erreur l'interpréteur explique 
brièvement, en anglais, la cause de l'erreur et désigne l'endroit du programme
où l'er\-reur a été trouvée (le numéro de la ligne, par exemple {\tt line 15}, puis la ligne 
elle-même et, en dessous, un ou plusieurs 
accents circonflexes {\idt \char94} sous le texte fautif). Le message commence par 
le nom de l'erreur, par exemple {\tt SyntaxError} (erreur de syntaxe). Il convient alors de
corriger l'erreur et de relancer le programme.

Une erreur découverte pendant l'exécution est signalée de la même façon, mais le 
message commence par {\tt Traceback (most recent call last):} (historique, l'appel le plus 
récent en dernier) et se termine par le nom de l'erreur, par exemple {\tt NameError} 
(erreur de nom).

\begin{exemple}[Programme syntaxiquement incorrect]{\rm
\hfill Le pro\-gram\-me {\idt CocotteErronee} de la page 
\pageref{prog:Cocotte-errone}
contient plusieurs des erreurs les plus fréquentes pour
un programme aussi simple.
Quand on lance ce programme {\idt CocotteErronee}, l'interpréteur commence par indiquer
`{\tt SyntaxError: unterminated triple-quoted string literal (detected at line 15)}' 
(chaîne entre triples guillemets non terminée, détectée à la ligne 15) en désignant 
le {\cha \char34\char34\char34} de la 
première ligne~: les triples guillemets fermants ont été oubliés à la fin du 
texte qui explique ce que fait le programme, si bien que tout le reste du 
programme, jusqu'à la dernière ligne 15, fait partie de ce texte$\:$!
Ayant ajouté les {\cha \char34\char34\char34} oubliés à la fin de la première ligne, 
une nouvelle tentative indique
`{\tt SyntaxError: \textquotesingle{}(\textquotesingle{} was never closed}' (la parenthèse `{\idt (}' n'a jamais été fermée) 
en désignant la parenthèse de {\idt st(} à la ligne 15. Nous la fermons en écrivant {\idt st()}.
Le programme est maintenant syntaxiquement correct et son exécution peut commencer. 
Elle s'arrête aussitôt, sur la première instruction, avec le message 
`{\tt NameError: name \textquotesingle{}phtd\textquotesingle{} is not defined}' (le nom {\idt phtd} n'est pas défini). 
Ayant oublié la ligne `{\mcl from}\ {\idt LibRobot}\ {\mcl import}\ {\idt *}', 
l'interpréteur ignore la signification des commandes du \robotErr, et la fenêtre 
du \robot\ ne s'est même pas ouverte. 
Après avoir ajouté cette ligne au début du programme, nous pouvons le relancer.
La fenêtre s'ouvre, le \robot\ dessine le premier trait de l'arrière de la cocotte, 
puis tout s'arrête avec le message 
`{\tt NameError: name \textquotesingle{}pgtg\textquotesingle{} is not defined. Did you mean: \textquotesingle{}phtg\textquotesingle{}?}' 
(le nom {\idt pgtg} n'est pas défini. Vouliez-vous dire {\idt phtg}$\:$?) qui désigne la 
ligne où se trouve {\idt pgtg()}. C'est une faute de
frappe que nous corrigeons en remplaçant {\idt  pgtg} par {\idt  phtg}. 
Nous obtenons ainsi laborieusement le programme {\idt Cocotte} de la page 
\pageref{prog:Cocotte} qui est correct (au commentaire près, écrit avec {\cmt \#} dans {\idt Cocotte}).
L'interpréteur s'arrête à la première erreur qu'il trouve~: il faut donc 
relancer le programme après chaque correction. On remarquera que les deux erreurs 
de syntaxe ont été signalées avant toute exécution, alors que les deux noms inconnus 
n'ont été découverts qu'au cours de l'exécution, au moment où l'or\-di\-na\-teur est arrivé 
à l'instruction fautive.
\noindent $\Box$ }\end{exemple}


\section{Exécution des programmes} 
%\label{sec:4:exec-progs} 
L'{\em exécution}\index{Exécution} d'un programme consiste à faire réaliser par
l'or\-di\-na\-teur les actions définies par les instructions du programme. Comme dans notre cas les instructions sont des
commandes du \robot, l'exécution du programme permet de piloter le
\robot. 

{\newcommand{\touche}[1]{{\begin{picture}(12,12)(0, 2)\put(0,0){\framebox(12,12)[lb]{}}\put(1.5,2){\makebox(0,0)[lb]{\tiny #1}}\end{picture}\hspace{0.3ex}}}
Quand le \robot\ a terminé son dessin, la commande {\idt st()} affiche en bas de 
la fenêtre les touches qui permettent de \hbox{continuer :}
\begin{list}{$\scriptscriptstyle\bullet$}{
\renewcommand{\rightmargin}{0mm}
}
\item diminuer (en tapant une ou plusieurs fois sur la touche moins \touche{$-$}) 
ou augmenter (en tapant une ou plusieurs fois sur la touche plus \touche{$+$}) 
la vitesse du \robot\ (qui peut varier de 0 à 10)$\:$;
\item faire redessiner tout le dessin par le \robot, à la vitesse choisie, en tapant 
sur une autre touche, par exemple la touche de retour à la ligne \framebox[11.mm][l]{$\hookleftarrow$~~~}\hspace{0.3ex} ou \hspace{0.3ex}\framebox[11.mm][l]{~~\tiny\texttt{return}}\hspace{0.3ex}$\:$; 
quand le \robot\ dessine trop rapidement pour pouvoir observer le
dé\-rou\-le\-ment du dessin, il suffit de diminuer sa vitesse puis de lui faire 
redessiner le dessin$\:$;
\item arrêter définitivement l'exécution du programme en tapant sur la 
touche d'échappement\index{Touche d'échappement} \touche{Esc} (ou sur l'une des touches 
\touche{\S} ou \touche{\#}), ce qui donne l'écran de la figure
\ref{fig:ecran-arret} (page \pageref{fig:ecran-arret}).
\end{list}

\index{Exécution (interruption)}\index{Interruption de l'exécution}Pendant 
l'exécution d'un programme de pilotage du \robot, nous
pouvons aussi arrêter le dessin du \robot\ à tout moment en tapant sur la touche 
d'échappement \touche{Esc} (ou sur \touche{\S} ou \touche{\#}). 
Le dessin reste alors affiché tel qu'il était à cet instant, jusqu'à ce que nous tapions 
sur une touche. Par exemple si nous interrompons l'exécution
du programme {\idt Stegosaure} de la page \pageref{prog:Stegosaure}
au moment où le \robot\ termine le dessin de la patte avant gauche,
nous obtenons un écran semblable à celui de la figure
\ref{fig:ecran-stop} de la page \pageref{fig:ecran-stop}.
}

\begin{figure}[p] %[htbp]
 \includevspacePostScript[7.6cm]{ecran-stop}{0.95}{0.95}
 \caption{Écran de l'or\-di\-na\-teur après un stop}\label{fig:ecran-stop}
\end{figure}

\begin{figure}[p] %[htbp]
 \includevspacePostScript[7.6cm]{ecran-arret}{0.95}{0.95}
 \caption{Écran de l'or\-di\-na\-teur avant l'arrêt définitif}
 \label{fig:ecran-arret}
\end{figure}

\clearpage
\section{Erreurs logiques} 
%\label{sec:4:err-exec} 

\index{Erreur logique}\index{Erreur de programmation}Après avoir éliminé toutes les erreurs syntaxiques c'est-à-dire 
les erreurs de programmation qui relèvent de la violation des règles de grammaire 
du langage \python, nous ne sommes pas encore au bout de nos peines avant de pouvoir
contempler le \hbox{dessin$\:$!} Il reste encore les
{\em erreurs logiques}
que nous découvrirons à l'exécution du programme.

La première erreur possible est que la succession de commandes que notre
programme demande à l'or\-di\-na\-teur de donner au \robot\ conduise celui-ci
en dehors du cadre. C'est le cas du programme suivant :
\input python/SortirDuCadre.tex
\noindent Ce programme donne le triste résultat 
ci-dessous, accompagné dans la fenêtre de terminal du message 
`{\tt Erreur ligne 8 de SortirDuCadre.py : le robot sortirait du cadre (av)}'~:
%indiqué à la figure \ref{fig:ecran-sortir} (page \pageref{fig:ecran-sortir}).
\begin{figure}[htbp]
 \includevspacePostScript[7.6cm]{ecran-sortir}{0.95}{0.95}
 \caption{Écran de l'or\-di\-na\-teur quand le robot sort du cadre}
 %\label{fig:ecran-sortir}
\end{figure}

La deuxième erreur possible est que nous nous soyons trompés dans
le programme et que le dessin obtenu ne ressemble pas du tout à ce
que nous avions imaginé. Par exemple, le programme {\idt StegosaureErrone}
ci-dessous est une version logiquement erronée mais syntaxiquement 
correcte du programme {\idt Stegosaure} de la page \pageref{prog:Stegosaure}. 
\input python/StegosaureErrone.tex
L'auteur de ce programme {\idt StegosaureErrone} a par deux fois confondu sa gauche 
et sa droite. Le 
résultat de l'exécution de ce programme tel
que représenté
ci-dessous
(figure \ref{fig:StegosaureErrone}) %page \pageref{fig:StegosaureErrone} 
éton\-ne\-rait plus d'un pa\-lé\-on\-to\-lo\-gue$\:$!
\begin{figure}[htbp]
 \includehvspacePostScript[105mm,6.8cm]{StegosaureErrone}{0.95}{0.95}
 \caption{Stégosaure plutôt raté}
 \label{fig:StegosaureErrone}
\end{figure}

Une troisième erreur, propre aux débutants en \python, est d'oublier les parenthèses 
{\idt ()} après une commande du \robot, par exemple `{\idt av;\ pqtd()}' au lieu de 
`{\idt av();\ pqtd()}'. Le programme est syntaxiquement correct mais le \robot\ n'avance pas 
car {\idt av} sans {\idt ()} désigne la commande sans demander de l'exécuter. 
L'interpréteur \python\ ne signale rien du tout~: ni erreur, ni même 
un {\em avertissement}\index{Avertissement}. 
Seul le dessin, auquel il manque un trait, révèle l'erreur. Il faut donc toujours 
vérifier que chaque commande est bien suivie de ses parenthèses.

\section{Comment localiser les erreurs logiques dans le programme} 
%\label{sec:4:trouv-err-exec} 

Pour trouver les erreurs logiques dans le programme il faut s'imaginer
que l'on exécute le programme à la place du \robotErr. Quand le programme 
est très long ceci n'est pas facile.
Le mieux est donc de suivre l'exécution du programme en mode {\em pas à pas},
\index{Exécution en pas à pas}
ce qui permet de suivre les commandes transmises par l'or\-di\-na\-teur au \robotErr\ pour localiser l'erreur avec précision. Quand l'erreur est
localisée, il faut arrêter l'exécution
du programme pour pouvoir la corriger dans le programme. Une fois
que c'est fait, on peut relancer le programme pour tenter une nouvelle
exécution.
{\newcommand{\touche}[1]{{\begin{picture}(12,12)(0, 2)\put(0,0){\framebox(12,12)[lb]{}}\put(1.5,2){\makebox(0,0)[lb]{\tiny #1}}\end{picture}\hspace{0.3ex}}}
\begin{list}{$\scriptscriptstyle\bullet$}{
\renewcommand{\rightmargin}{0mm}
}
\item Pour passer en mode pas à pas nous \hbox{pouvons :}
\end{list}
\begin{list}{-}{
\renewcommand{\labelwidth}{3mm}
\renewcommand{\labelsep}{3mm}
\renewcommand{\leftmargin}{15mm}
\renewcommand{\rightmargin}{0mm}
}
  \item à la fin du dessin, réduire la vitesse à zéro en tapant plusieurs fois sur la touche
  \touche{$-$} puis faire redessiner le dessin en tapant sur la touche de retour à la ligne
  \framebox[11.mm][l]{$\hookleftarrow$~~~}\hspace{0.3ex} ou \hspace{0.3ex}\framebox[11.mm][l]{~~\tiny\texttt{return}}$\:$;
  \item ou insérer une instruction {\idt pp()} dans le programme \python, un
  peu avant l'endroit supposé de l'erreur. Cette instruction {\idt pp()} fait passer
  l'or\-di\-na\-teur en mode d'exécution {\underline p}as à {\underline p}as.
\end{list}
\begin{list}{$\scriptscriptstyle\bullet$}{
\renewcommand{\rightmargin}{0mm}
}
\item En mode pas à pas, l'or\-di\-na\-teur attend que nous lui donnions
l'ordre d'exécuter chaque commande du \robot\ en tapant sur une touche, par exemple la 
touche de retour à la ligne \framebox[11.mm][l]{$\hookleftarrow$~~~}\hspace{0.3ex} ou \hspace{0.3ex}\framebox[11.mm][l]{~~\tiny\texttt{return}}. On peut donc suivre
le déroulement de l'exécution du programme instruction après
instruction, ce qui permet en général de trouver assez facilement l'erreur.
\enlargethispage{2\baselineskip}% touches return ajoutees page 50 : garder la fin de la liste sur la page
\item Quand l'erreur est localisée avec précision dans le programme, on
 peut arrêter l'exécution du programme en tapant sur
 la touche d'é\-chap\-pe\-ment \touche{Esc} (ou sur \touche{\S} ou \touche{\#}).
\item[]  Il est également possible de programmer le retour à une vitesse
   d'\underline{ex}é\-cu\-tion normale en insérant une instruction {\idt ex()}
  dans le programme.
\end{list}
}
Il faut presque toujours plusieurs essais avant que le programme soit 
lo\-gi\-que\-ment correct.
Quand il l'est, on est souvent tenté de l'améliorer. Le plus difficile
est alors de savoir quand \hbox{s'arrêter$\:$!}

\begin{table}[htbp]
\begin{center}
\begin{tabular}{@{}|l|p{11cm}|@{}} \hline
{\idt st()} & \underline{st}oppe l'exécution du programme à la fin du dessin. 
             La vitesse peut être diminuée par
 \begin{picture}(12,12)(0, 2)\put(0,0){\framebox(12,12)[lb]{}}\put(1.5,2){\makebox(0,0)[lb]{$\scriptscriptstyle -$}}\end{picture}\hspace{0.3ex} (minimum 0) ou augmentée par
 \begin{picture}(12,12)(0, 2)\put(0,0){\framebox(12,12)[lb]{}}\put(1.5,2){\makebox(0,0)[lb]{$\scriptscriptstyle +$}}\end{picture}\hspace{0.3ex} (maximum 10).
			 Toute autre touche fait redessiner le dessin à la vitesse indiquée. 
			 L'exécution  s'arrête définitivement en tapant sur 
			 \begin{picture}(12,12)(0, 2)\put(0,0){\framebox(12,12)[lb]{}}\put(1.5,2){\makebox(0,0)[lb]{\tiny esc}}\end{picture}\hspace{0.3ex} ou 
  \begin{picture}(12,12)(0, 2)\put(0,0){\framebox(12,12)[lb]{}}\put(1.5,2){\makebox(0,0)[lb]{\tiny \S}}\end{picture}\hspace{0.3ex} ou
  \begin{picture}(12,12)(0, 2)\put(0,0){\framebox(12,12)[lb]{}}\put(1.5,2){\makebox(0,0)[lb]{\tiny \#}}\end{picture}\hspace{0.3ex}.\\ \hline
{\idt pp()} & exécute la suite du programme en mode \underline{p}as
             à \underline{p}as (vitesse 0)~: chaque commande du \robot\ 
			s'exécute en tapant sur une touche.\\ \hline
{\idt ex()} & exécute la suite du programme en mode \underline{ex}écution
             normale (vitesse 10). \\ \hline
\end{tabular}
\end{center}
\label{tab:cmd-vit1}
\caption{Commandes de contrôle du mode de pilotage du \protect\robot}
\end{table}

\noindent \rule{137mm}{0.1mm} %---------------------------------------------
\begin{exercices}{\ \\

\vspace{-\baselineskip}
\noindent \begin{minipage}[t]{88mm}
\tiret Écrire, mettre au point et exécuter des programmes \python\ de dessin du
\robot\ tel qu'il est stylisé ci-contre avec son
crayon baissé et son crayon levé.
\end{minipage}
\hfill
\begin{minipage}[t]{44mm}
\vspace{2.0cm}\includePostScript{Robot+BC+LC}{0.95}{0.95}
\hspace{-4.4cm}  % 4.4cm de large
\end{minipage}

\noindent \begin{minipage}[t]{103mm}
\tiret  Écrire, mettre au point et exécuter un programme \python\ qui
pilote le \robot\ pour lui faire dessiner le chien représenté
ci-contre.
\end{minipage}
\hfill
\begin{minipage}[t]{29mm}
\vspace{2.3cm}\includePostScript{Chien}{0.95}{0.95}
\hspace{-2.9cm}  % 2.9cm de large
\end{minipage}

\noindent \begin{minipage}[t]{103mm}
\tiret  Programmer les dessins de la
pa\-ge \pageref{fig:quadrilateres}.
\end{minipage}

\vspace{1.5mm}

\tiret  Programmer le dessin de vos initiales en lettres majuscules
dans le style de la lettre {\tt H} en relief représentée à la
figure \ref{fig:h-relief} (page \pageref{fig:h-relief}).

\noindent $\Box$}\end{exercices}

\newpage
\begin{corrige}{Nous donnons un programme, les autres se trouvant dans le dossier {\idt livre}.
\input python/Chien.tex
\noindent $\Box$}\end{corrige}

\vfill

%%%%%%%%%%%%%%%%%%%%%%%%%%%%%%%%%%%%%%%%%%%%%%%%%%%%%%%%%%%%%%%%%%%%%%%%%%%
\noindent \framebox[\textwidth]{
\begin{minipage}{\myparboxwidth}
\setlength{\parindent}{2em}
\small \sf
\vspace{\abovedisplayskip}
\section*{Addenda}
Le cycle écriture, mise au point syntaxique, traduction, mise au point
logique et exécution des programmes s'est considérablement
raccourci ces dernières années grâce aux progrès fantastiques
de l'architecture des or\-di\-na\-teurs qu'il convient de comprendre très
schématiquement. 
Un micro-or\-di\-na\-teur \hbox{comprend :}
\begin{list}{$\scriptscriptstyle\bullet$}{
}
\item une unité arithmétique et logique pour faire des calculs et
manipuler des données,
\item une mémoire pour conserver temporairement des données
et des programmes (son contenu est perdu quand on coupe l'électricité),
\end{list}

\vspace{\abovedisplayskip}
\end{minipage}
}
%%%%%%%%%%%%%%%%%%%%%%%%%%%%%%%%%%%%%%%%%%%%%%%%%%%%%%%%%%%%%%%%%%%%%%%%%%%

%%%%%%%%%%%%%%%%%%%%%%%%%%%%%%%%%%%%%%%%%%%%%%%%%%%%%%%%%%%%%%%%%%%%%%%%%%%
\noindent \framebox[\textwidth]{
\begin{minipage}{\myparboxwidth}
\setlength{\parindent}{2em}
\small \sf
\vspace{\abovedisplayskip}
\begin{list}{$\scriptscriptstyle\bullet$}{
}
\item un ou plusieurs disques durs magnétiques ou à état solid (SSD)
 pour conserver en permanence des fichiers contenant des données ou des programmes,
\item des prises (USB) pour brancher des clés de mémoire qui permettent de
conserver en permanence des fichiers
contenant des données ou des programmes ou les échanger avec d'autres
or\-di\-na\-teurs,
\item un clavier pour entrer des données ou des programmes en mémoire,
\item un écran graphique pour visualiser des textes et des images,
\item une souris pour désigner des points sur l'écran graphique,
\item un bus pour permettre des échanges rapides de données et de
programmes entre ces différents organes.
\end{list}

Quand on met l'or\-di\-na\-teur en marche, il est cablé pour charger dans
sa mémoire et exécuter un programme rangé en permanence sur disque
que l'on appelle le {\em système d'exploitation}\index{Système d'exploitation}. Ce programme a pour
tâche essentielle de gérer tous les organes de l'or\-di\-na\-teur et de
dialoguer avec l'utilisateur grâce à un {\em interpréteur de commandes}\index{Interpréteur de commandes}
qui va attendre que l'utilisateur lui donne des ordres par l'intermédiaire
du clavier et les exécuter immédiatement les uns après les autres
jusqu'à ce que l'utilisateur lui donne l'ordre d'éteindre
l'or\-di\-na\-teur. Parmi les commandes de l'interpréteur, citons par exemple la
suppression d'un fichier rangé sur disque ou le lancement de
l'exécution d'un programme écrit en langage machine (généralement
appelé {\em programme d'application}\index{Programme d'application}) comme par exemple un {\em
éditeur de texte} pour créer un texte et le ranger dans un fichier
ou un {\em interpréteur} pour traduire un programme \python\ en langage machine et le faire 
exécuter.

Les premières interfaces entre l'homme et la machine étaient purement textuelles. On tape sur
le clavier et on voit le résultat sur une imprimante vite remplacée par un écran avec seulement
des lignes de texte de longueur maximale 80 caractères. Les écrans en noir et blanc furent
remplacés pat des écrans couleur (comme on voit encore dans l'application \texttt{terminal}).
Puis sont apparus les écrans graphiques constitués de très petits pixels de couleur. Le
fenêtrage permet de donner l'illusion de plusieurs écrans graphiques sur un seul. Puis vint
l'idée d'icônes pour représenter les fichiers et des applications sur lesquels on peut
double cliquer pour lancer une application. Par exemple,  un double click sur l'icône de
\texttt{LeRobotPython.app} lance le programme de dessin interactif tandis que faire glisser l'icône
d'un programme du \robot\ sur cette icône de l'application \texttt{LeRobotPython.app} exécute ce programme
en faisant apparaître la fenêtre de dessin et le terminal. Noter qu'il faut cliquer sur
l'une ou l'autre des fenêtres pour l'utiliser. On peut arrêter le programme par Control+C ou
ctrl+C (que les informaticiens notent \^{}C) en sélectionnant la fenêtre terminal ou en fermant cette fenêtre ou encore avec la touche  d'échappement (\texttt{esc}) dans la fenêtre graphique.

% (pas d'espace en bas de ce cadre : il tient dans la page)
\end{minipage}
}
%%%%%%%%%%%%%%%%%%%%%%%%%%%%%%%%%%%%%%%%%%%%%%%%%%%%%%%%%%%%%%%%%%%%%%%%%%%

%%%%%%%%%%%%%%%%%%%%%%%%%%%%%%%%%%%%%%%%%%%%%%%%%%%%%%%%%%%%%%%%%%%%%%%%%%%
\noindent \framebox[\textwidth]{
\begin{minipage}{\myparboxwidth}
\setlength{\parindent}{2em}
\small \sf
\vspace{\abovedisplayskip}

Les or\-di\-na\-teurs sont {\em binaires} en ce sens que toutes les données, en
particulier les programmes, sont représentés par des suites de 0 et
de 1. L'or\-di\-na\-teur fonctionne en transférant une instruction du programme
depuis une case de sa mémoire
dans son unité arithmétique et logique (via le bus) et en 
l'exécutant (ce qui en général modifie le contenu de la mémoire,
du disque, de l'écran, etc.) puis en recommençant ainsi avec l'instruction
rangée dans la case suivante de sa mémoire. Certaines instructions
dites {\em instructions de branchement} permettent d'indiquer que l'instruction suivante
est rangée dans une case mémoire de numéro donné.


Pour exécuter un programme écrit en langage machine rangé sur disque,
le système d'exploitation demande à l'unité arithmétique et
logique de charger ce programme en mémoire c'est-à-dire de
recopier le programme dans un endroit de la mémoire
qui était inoccupé puis d'exécuter une instruction de branchement
à la première instruction de ce programme. Le programme doit se terminer
par une instruction de branchement en retour vers le système 
d'exploitation. Les or\-di\-na\-teurs modernes possèdent également la fa\-cul\-té
d'interrompre l'exécution du programme en cours puis de la reprendre
ultérieurement comme si rien ne s'était passé. Entre temps, c'est
le système d'exploitation qui s'exécute pour réagir en fonction de
la cause de l'interruption. Par exemple si l'utilisateur a bougé la souris,
ses nouvelles coordonnées sont rangées en mémoire. Si l'horloge
sonne, ce qui arrive tous les soixantièmes de seconde par exemple, l'heure 
est mise à jour, etc. En particulier l'utilisateur dispose
de boutons ou de touches (%touche Esc (touche d'echappement)
\begin{picture}(14,12)(0, 2) 
\put(0,0){\framebox(14,12)[lb]{}}
\put(1.5,2){\makebox(0,0)[lb]{\tiny Esc}}
\end{picture}\hspace{0.3ex},
%touche §
\begin{picture}(12,12)(0, 2)
\put(0,0){\framebox(12,12)[lb]{}}
\put(1.5,2){\makebox(0,0)[lb]{\tiny \S}}
\end{picture}\nolinebreak\hspace{0.3ex}\nolinebreak,
%touche #
\begin{picture}(12,12)(0, 2)
\put(0,0){\framebox(12,12)[lb]{}}
\put(1.5,2){\makebox(0,0)[lb]{\tiny \#}}
\end{picture}\hspace{0.3ex} 
 pour notre \robot) permettant d'interrompre
l'exécution de son programme s'il constate une erreur.


Parmi les applications que peut exécuter le système d'exploitation, 
il y en a une qui nous intéresse particulièrement, c'est l'interpréteur \python.

Les environnements de programmation modernes (comme {\tt IDLE}, fourni avec \python) 
se chargent de l'édition du texte du
programme, c'est-à-dire de sa création et de ses modifications
par interaction avec l'utilisateur. Le 
programme \python\ est affiché sur l'écran et rangé en mémoire 
au fur et à mesure que
l'utilisateur le tape au clavier. Il est prévu qu'il puisse le modifier
en utilisant diverses touches de fonctions et la souris. La possibilité
lui est également offerte de sauvegarder son texte sur le disque dur
dans un fichier dont il doit donner le nom.

\vspace{\abovedisplayskip}
\end{minipage}
}
%%%%%%%%%%%%%%%%%%%%%%%%%%%%%%%%%%%%%%%%%%%%%%%%%%%%%%%%%%%%%%%%%%%%%%%%%%%

%%%%%%%%%%%%%%%%%%%%%%%%%%%%%%%%%%%%%%%%%%%%%%%%%%%%%%%%%%%%%%%%%%%%%%%%%%%
\noindent \framebox[\textwidth]{
\begin{minipage}{\myparboxwidth}
\setlength{\parindent}{2em}
\small \sf
\vspace{\abovedisplayskip}

Quand l'édition de texte est terminée, l'environnement peut
faire traduire le programme en un programme équivalent
écrit en langage machine. Pour les langages dits {\em compilés}\index{Compilation}, 
comme \ocaml\ ou C, cette traduction est faite une fois pour toutes par 
un {\em compilateur}\index{Compilateur}~: le programme compilé est rangé
sur disque, puis l'utilisateur retourne au niveau du système d'exploitation pour 
charger en mémoire
et exécuter son programme compilé comme expliqué précédemment.

Pendant la phase de débogage du programme, il est plus rapide de 
ranger le programme traduit directement en mémoire, le traducteur
se chargeant lui-même de son exécution puis reprenant la main quand
cette exécution est terminée ou interrompue. C'est ainsi que procède 
l'interpréteur \python, à chaque fois qu'on lui demande d'exécuter un programme.

\vfill

On comprend les progrès considérables qui ont été faits depuis les
premiers or\-di\-na\-teurs pour lesquels il fallait modifier le cablage pour
changer le programme \cite{levy1} :
\begin{list}{$\scriptscriptstyle\bullet$}{
}
\item le programme étant rangé en mémoire (selon l'idée de
Von Neumann\footnote{\ \href{https://en.wikipedia.org/wiki/John_von_Neumann}{John von Neumann}, 1903--1957, mathématicien américain.})
et la mémoire reliée à des disques rapides, il est possible de
changer le programme de la machine presque \hbox{instantanément$\:$;}
\item l'usage de langages de haut niveau comme \python\ ou \ocaml\ et d'interpréteurs 
ou de compilateurs
performants, permet d'éviter d'avoir à programmer en un langage
machine \hbox{ingrat$\:$;}
\item l'écran alphanumérique permet à l'utilisateur d'éditer son programme
de manière interactive au clavier, ce qui fait gagner un temps 
considérable par rapport aux manipulations de cartes perforées en carton
encore utilisées dans les années \hbox{1970$\:$;}
\item la puissance de calcul, liée à la densité des transitors sur
les puces informatiques (qui varie de quelques milliards à plusieurs centaines de milliards
de transistors en 2016) permet de gérer des écrans
graphiques directement reliés à la mémoire, ce qui permet des
modifications très rapides et donc l'usage de fenêtres, souris, \hbox{etc$\:$;}
\item la taille de la mémoire (quelques milliards de caractères, Go, Gigaoctets) et
des disques (quelques milliers de milliards de caractères, To, Teraoctets) permet d'envisager des
programmes et des données de plus en plus grands et donc des applications
de plus en plus sophistiquées.
\end{list}
\vspace{\abovedisplayskip}
\end{minipage}
}
%%%%%%%%%%%%%%%%%%%%%%%%%%%%%%%%%%%%%%%%%%%%%%%%%%%%%%%%%%%%%%%%%%%%%%%%%%%

%%%%%%%%%%%%%%%%%%%%%%%%%%%%%%%%%%%%%%%%%%%%%%%%%%%%%%%%%%%%%%%%%%%%%%%%%%%
\noindent \framebox[\textwidth]{
\begin{minipage}{\myparboxwidth}
\setlength{\parindent}{2em}
\small \sf
\vspace{\abovedisplayskip}
La taille des plus grands programmes que l'on peut actuellement écrire
se chiffre maintenant en dizaines ou centaines de millions de lignes, ce qui est considérable
et ne va pas sans difficultés car on atteint les limites de la
compréhension humaine. Une formation de l'esprit, une 
nouvelle rigueur, plus précise encore que la logique des mathématiciens,
devient la qualité première qui, en plus de l'imagination et du sens de
l'esthétique, fait les bons informaticiens.

\indent Un peu de méthode aide
également beaucoup. Toute l'histoire des langages de programmations
est remplie d'inventions permettant de formuler plus clairement la façon
d'organiser les données en mémoire ou les calculs de l'or\-di\-na\-teur. A
chaque fois les capacités humaines, donc la taille limite des plus grands
 programmes, ont été repoussées. 
Une idée très ancienne de structuration des calculs, due à Von Neumann et Turing%
\footnote{\ \href{https://en.wikipedia.org/wiki/Alan_Turing}{Alan Turing}, 1912--1954, mathématicien anglais.} est celle
de {\em procédure}, on disait autrefois {\em routine}
ou {\em sous-programme}. Nous commençons à l'introduire dans le chapitre suivant.
\vspace{\abovedisplayskip}
\end{minipage}
}
%%%%%%%%%%%%%%%%%%%%%%%%%%%%%%%%%%%%%%%%%%%%%%%%%%%%%%%%%%%%%%%%%%%%%%%%%%%
\newpage
\thispagestyle{empty}